\documentclass{beamer}
% \documentclass[handout]{beamer}
\input{inputs/preamble}

\title
{A Presentation Template}

\subtitle
{With a Subtitle, if Needed}

\date
{March 2020}

\begin{document}

\begin{frame}[plain]
  \titlepage
\end{frame}

\begin{frame}{On Science}
  \vfill
  \begin{itemize}
    \item<1-> Science is \emph{cool!} \citep{Sagan}
      \vfill
    \item<2-> Unfortunately \emph{a big great many lot} of people do not seem to get it\dots \citep{Feynman74}
      \vfill
    \item<3-> And now some \emph{math:}
      \vfill
      \begin{theorem}[Euler's identity] % Note the ( and ) in the optional name
        \begin{equation}
          e^{i\pi}+1=0
        \end{equation}
      \end{theorem}
  \end{itemize}
  \vfill
\end{frame}

\begin{frame}[fragile]{Verbathings}
  \begin{itemize}
    \item Not verbatim text, vs.\ \verb+verbatim text+
    \item The \verb+verbatim+ \emph{environment} also works:
      \begin{verbatim}
Lorem ipsum dolor sit amet, blah blah...
      \end{verbatim}
  \end{itemize}
\end{frame}

\begin{frame}{Summary}
  \begin{itemize}
  \item
    The \alert{first main message} of your talk in one or two lines.
  \item
    The \alert{second main message} of your talk in one or two lines.
  \item
    Perhaps a \alert{third message}, but not more than that.
  \end{itemize}
\end{frame}

\begin{frame}{References \& Notes}
  \vfill
  \bibliography{sources}
  \bibliographystyle{plainnat}
  \vfill
  \vfill
  \textbf{But beware:} \texttt{beamer} does not always play nice with references!\footnote{In particular, with overly fancy bibliographic styles\dots}\\[1em]
  Besides, in a presentation the number of references should be minimised anyway\dots
  \vfill
\end{frame}

\begin{frame}[plain]{Thank you for listening!}
  \begin{textblock*}{0.1\textwidth}(0.2\textwidth, 0.25\textheight)
    \LARGE \textit{Questions}
  \end{textblock*}
  \begin{centering}
    \hfill\hfill
		\begin{tikzpicture}[ultra thick,scale=0.7]
			\draw[line width=3.5,line cap=round] (1.5,0) .. controls ++(0,4) and 
			++(0,-2) .. (4,4)
			to[out=90,in=0] (2,6)
			to[out=180,in=90] 
			(0,4);
    \end{tikzpicture}
    \hfill
  \end{centering}

  \begin{textblock*}{0.1\textwidth}(0.55\textwidth, 0.02\textheight)
    \Huge \Smiley
  \end{textblock*}
\end{frame}

\end{document}

% vim: spell spelllang=en
